\section{Related Works}

\begin{table*}[h]
\centering
\newcommand{\cyes}{{\color{green!60!black}\checkmark}}
\newcommand{\cno}{{\color{red}\xmark}}
\resizebox{\textwidth}{!}{%
\begin{tabular}{lcccccccccc}
\toprule
\textbf{\multirow{2}{*}{Benchmark}}                              & \textbf{\multirow{2}{*}{Focus}} & \textbf{Num.}   & \textbf{Num.}   & \textbf{Num.}   & \textbf{Avg.}   & \textbf{DB}    & \textbf{Avg.} & \textbf{Refusal}  & \textbf{Human Task} & \textbf{Human} \\
                                                                 &                     & \textbf{Domains} & \textbf{Tasks}  & \textbf{Tools}  & \textbf{Steps}  & \textbf{Tables} & \textbf{FK}     & \textbf{Ability?} & \textbf{Curation?}  & \textbf{Plans?} \\
\midrule
\textit{General Tool Use}                                        &                     &                 &                 &                 &                 &                 &                       &                   &                     & \\
API-Bank \citep{li-etal-2023-api}                                & Tool-use             & 8          & 314             & 73              & ~3            & 0                & 0                     & \cno            & \cyes               & \cno \\
ACEBench \citep{chen2025acebenchwinsmatchpoint}                                  & Tool-use             & 8         & 2000            & 4538            & ~2           & 0                & 0                     & \cyes            & \cyes               & \cno \\
$\tau$-bench \citep{yao2024taubenchbenchmarktoolagentuserinteraction}             & User Interaction     & 2         & 165             & 28               & \textemdash             & 3                & 0.7                   & \cyes        & \cyes               & \cno \\
$\tau^2$-bench \citep{barres2025tau2}                                        & User Interaction       & 3                & 279             & 56               & \textemdash             & 9                & \textemdash                   & \cyes        & \cno               & \cno \\
\midrule
\textit{Enterprise Specific}                                     &                     &                 &                 &                 &                 &                 &                       &                   &                     & \\
WorkArena \citep{workarena2024}                                  & ServiceNow          & 7                & 33              & 30             & ~10             & 7                & 0.9                   & \cno            & \cno               & \cno \\
WorkArena++ \citep{boisvert2024workarenacompositionalplanningreasoningbased}     & ServiceNow          & 7                & 682 ($^{\dagger}$341)            & 30             & 30-50           & 7                & \textemdash                   & \cyes        & \cno               & \cno \\
ITBench \citep{jha2025itbenchevaluatingaiagents}                                 & IT                  & 3                & 94              & 10              & \textemdash             & \textemdash              & \textemdash                   & \cno            & \cyes          & \cno \\
WorkBench \citep{styles2024workbenchbenchmarkdatasetagents}                      & Workplace           & 5                & 690 ($^{\dagger}$69)             & 26              & 2             & 5                & 0                     & \cno             & \cyes          & \cno \\
TheAgentCompany \citep{xu2024theagentcompanybenchmarkingllmagents}               & Startup             & 7                & 175             & \textemdash             & \textemdash             & 0              & 0                     & \cno            & \cyes          & \cno \\
CRMArena \citep{huang-etal-2025-crmarena}                                       & Salesforce          & 1                & 1170 ($^{\dagger}$9)             & 27              & \textemdash             & 16               & 1.3                   & \cno            & \cyes          & \cno \\
CRMArena-Pro \citep{huang-etal-2025-crmarena-pro}                                & Salesforce          & 3                & 8560 ($^{\dagger}$19)            & \textemdash             & \textemdash             & 25               & \textemdash                   & \cyes        & \cyes          & \cno \\
EnterpriseBench \citep{vishwakarma2025llmshelpworksandbox}                       & Enterprise          & 5                & 500             & 46              & 3               & 17               & 1.2                   & \cyes        & \cno               & \cno \\ 
\midrule
\rowcolor{gray!15}
\textbf{\ours{} (Ours)}                                          & \textbf{Enterprise} & \textbf{8}       & \textbf{1150}   & \textbf{512}    & \textbf{9$^*$}     & \textbf{164}     & \textbf{1.7}          & \textbf{\cyes}   & \textbf{\cyes}     & \textbf{\cyes} \\
\bottomrule
\end{tabular}%
}
\caption{\textbf{Comparison with existing agentic benchmarks.} \textit{DB Tables} reports the number of unique database tables in the environment; \textit{Avg. FK} measures average foreign keys per table, indicating relational density and the complexity of inter-table dependencies agents must navigate. \textemdash{} denotes values not reported in the original work. $^*$Avg. Steps reflects ideal human-authored execution trajectories; model trajectories may require significantly more steps. Human Plans refer to step-by-step natural language plans written by experts to complete the task. $^{\dagger}$Parenthetical values indicate the number of unique task templates.}
\label{tab:related_works_comparison}
\end{table*}

LLM agent benchmarks broadly fall into three thematic groups: general-purpose tool-use and API evaluation, enterprise platform simulation, and agentic planning and computer-use. \Cref{tab:related_works_comparison} summarizes how \dataset{} compares across key dimensions.

\paragraph{API and Tool-Use Benchmarks}
Early benchmarking efforts focused on testing agents' ability to call APIs and chain tools in general, open-domain settings. ToolLLM \citep{qin2024toolllm} establishes large-scale evaluation across over 16,000 real-world web APIs, while API-Bank \citep{li-etal-2023-api} provides a smaller, runnable system for assessing planning and retrieval across 73 tools. ACEBench \citep{chen2025acebenchwinsmatchpoint} scales this further with 4,538 tools and dynamic multi-turn evaluation, and $\tau$-bench \citep{yao2024taubenchbenchmarktoolagentuserinteraction} and $\tau^2$-bench \citep{barres2025tau2} introduce user simulation and refusal robustness as evaluation axes. These benchmarks make important contributions to measuring tool-calling accuracy but remain anchored to general web-oriented or open-domain APIs. They do not model the multi-system, policy-constrained, stateful tool ecosystems that characterize enterprise environments, which is the setting \ours{} specifically targets.

\paragraph{Enterprise Benchmarks}
A second class of benchmarks targets enterprise platforms directly. Single-platform environments restrict scope to one vendor: WorkArena \citep{workarena2024} and WorkArena++ \citep{boisvert2024workarenacompositionalplanningreasoningbased} evaluate compositional task completion within ServiceNow, CRMArena \citep{huang-etal-2025-crmarena} and CRMArena-Pro \citep{huang-etal-2025-crmarena-pro} focus on Salesforce specific workflows, and ITBench \citep{jha2025itbenchevaluatingaiagents} targets IT incident resolution. Multi-domain efforts broaden the scope but with different emphases: TheAgentCompany \citep{xu2024theagentcompanybenchmarkingllmagents} simulates a startup software company requiring agents to interact via web browser, bash terminal, and code execution, a paradigm fundamentally different from the structured tool-calling API setting of \ours{}; WorkBench \citep{styles2024workbenchbenchmarkdatasetagents} covers office productivity tools; and EnterpriseBench \citep{vishwakarma2025llmshelpworksandbox} evaluates function-calling on sandboxed enterprise data. Across these benchmarks, relational database complexity is consistently limited to fewer than 25 tables, and tasks are either confined to a single vendor ecosystem or lack expert-authored grounding or explicit evaluation of refusal behavior. \ours{} addresses these limitations by spanning eight enterprise domains across both operational systems (CSM, ITSM, HR) and collaboration services (Email, Calendar, Teams, Drive), with 164 database tables with high connectivity, 512 tools, 1,150 expert-curated tasks including policy-constrained infeasible scenarios.

\paragraph{Agentic Planning and Computer-Use}
Planning is a core capability for autonomous agents, and several benchmarks have studied it across diverse general-purpose settings. UserBench \citep{qian2025userbenchinteractivegymenvironment} evaluates agents on multi-turn travel planning tasks where simulated users express preferences incrementally and implicitly, requiring proactive intent elicitation. Gaia2 \citep{froger2025arescalingagentenvironments} evaluates agents in an asynchronous simulated mobile environment, testing search, execution, temporal reasoning, adaptability, ambiguity handling, and multi-agent collaboration. VitaBench \citep{he2025vitabenchbenchmarkingllmagents} evaluates agents on complex multi-turn tasks drawn from real-world services like food delivery, in-store dining, and online travel using a simulated user with dynamic preferences. A parallel line of work focuses on computer-use agents: OSWorld \citep{xie2024osworld}, WindowsAgentArena \citep{bonatti2024arena}, WebArena \citep{zhou2024webarenarealisticwebenvironment} and UI-Vision \citep{nayak2025uivisiondesktopcentricguibenchmark} benchmark agents on desktop and web GUIs requiring complex sequential decision-making. While these settings demand sophisticated planning, they operate in everyday scenarios or computing environments rather than the policy-governed, multi-system contexts of enterprise environments. \ours{} shares the emphasis on extended planning horizons but is uniquely positioned at the intersection of planning depth and enterprise domain fidelity.
